\section{Conclusion}

This work presented TAMR-DTI, a multimodal DTI framework that represents drug-target pairs through
target-aware conformer modulation, ligand-conditioned protein encoding, protein-side Mamba refinement,
and BiIntention-based interaction matching. Across five repeated runs, the model achieves the strongest
reported AUROC/AUPRC values on all four benchmark datasets considered in this study: 0.964/0.954 on
BindingDB, 0.930/0.935 on BioSNAP, 0.982/0.978 on Human, and 0.992/0.991 on C. elegans. These results
show that the proposed representation path improves the primary ranking objective for DTI prediction across
datasets with different sizes, label ratios, and organism sources.

The diagnostic results further clarify where the gains come from. On BioSNAP, removing target-aware
conformer aggregation lowers AUROC from 0.933 to 0.916, while removing ligand-conditioned protein FiLM
or protein-side Mamba refinement also reduces both ranking and fixed-threshold performance. The conformer-
count comparison shows that the eight-conformer target-aware variant reaches the best ranking scores among
the tested geometry settings (0.9322 AUROC and 0.9356 AUPRC), and the conformer-weight analysis shows
non-uniform, sample-specific conformer usage over the held-out test pairs. Together, these findings support
the central design conclusion: ligand geometry should be treated as pair-level evidence conditioned on the
protein context, while long-range protein refinement should complement rather than replace explicit drug-
target matching.

The broader significance is that DTI models can move beyond independently encoded drug and protein
descriptors toward pair-specific representation construction. More concretely, TAMR-DTI provides a
practical way to select and aggregate generated ligand conformers according to the target, feed the resulting
ligand context back into protein representation learning, refine the ordered protein sequence with an
efficient state-space block, and retain an explicit bidirectional interaction module for prediction. This gives
future DTI systems a modular template for combining molecular geometry, protein sequence context, and
interpretable pairwise matching. In practical drug-discovery workflows, such a model can help prioritize
compound-target pairs for downstream validation, narrow the candidate space before costly biochemical
assays, and provide model-internal conformer-usage evidence for why a ligand may be more plausible for one
target than another, while still requiring experimental validation before biological or clinical use. The
current evidence is strongest for ranking performance; future work should strengthen the conclusion with
matched multi-seed ablations, harder generalization splits, and more direct structural supervision for conformer
selection.
